\section{Introduction}%
\label{sec:introduction}

While deep learning has become popular over the last decade, its predictions are hardly explainable to humans, and attempts of peeking into the ``black box'' often focus on \emph{local} explainability rather than \emph{global} explainability, i.e., on single predictions rather than the model as a whole~\citep{Guidotti2019Survey}. However, when classifying nodes in knowledge graphs, e.g., predicting the type of entities, predicting which molecules cause cancer, or predicting which drugs are promising treatment candidates, it is often desirable to understand the model as a whole. This enables experts to double-check the model's plausibility and adapt the learned rules to reflect domain knowledge. Toward this end, interpretable models based on description logics have been proposed~\citep{Lehmann2009DL-Learner,Tran2017Parallel,Westphal2019SML-Bench}. Concepts in description logics can be directly mapped to OWL class expressions allowing to make use of the rich data ecosystem centered around the Web Ontology Language, OWL, the W3C standard at the core of the Semantic Web~\cite{Li2021Computing, Hogan2021Knowledge}. Tools for verbalizing OWL class expressions in natural language are readily available to explain the expressions to users~\citep{Moussallem2020NABU, Moussallem2020Generating, Ngonga2019Holistic}. OWL class expressions can be applied to web-based knowledge graphs such as DBpedia~\cite{Auer2007DBpedia}, Wikidata~\cite{Vrandecic2014Wikidata}, and YAGO~\cite{Suchanek2007Yago}, as well as to semantic data on websites, e.g., in the schema.org ontology~\cite{Hernich2015Schema.org} for type prediction~\cite[cf.][]{Paulheim2013Type, Zahera2021ASSET}, product classification~\cite{Zhang2019Product}, and discovery and querying of web APIs~\cite{Michel2019Enabling,Wolters2017Linking}. 

State-of-the-art approaches for concept learning are based on inductive logic programming with refinement operators~\citep{Fanizzi2018DLFoil,Lehmann2009DL-Learner,Lehmann2010Concept}: Starting from the most general concept $\top$ (also known as \texttt{Thing}), they refine the concept iteratively. For example, a concept is replaced by a subconcept or by multiple subconcepts joined via logical operators such as disjunction, conjunction or negation. In these approaches, the generation of candidate concepts is almost exclusively based on the ontology, and positive and negatives examples are only used to evaluate the generated concepts. In case of large ontologies, these approaches lead to a combinatorial explosion and many generated concepts never appear in the instance data leading to an unnecessarily large search space. For example, in the DBpedia ontology, a person can have over~100 different properties and ILP approaches would try all of them including combinations thereof. However, most of these properties (and combinations thereof) rarely occur in the instance data, e.g., it is seldom stated whether a person is left-handed or right-handed (object property ``handedness'') and neither is their hip size (data property ``hip size'').

Moreover, many knowledge bases contain large amounts of data properties connecting entities to literals, such as numeric values, which are crucial for good predictions. However, most state-of-the-art approaches for concept learning neglect data properties. Only \mbox{DL-Learner} \citep{Lehmann2010Concept} and \mbox{SPaCEL}~\citep{Tran2017Parallel} have some rudimentary support: DL-Learner divides the value range of a data property into bins containing approximately the same number of examples; \mbox{SPaCEL} considers all thresholds in a brute-force manner. Neither approach takes the distribution of positive and negative examples into account, resulting in suboptimal thresholds or long runtimes. DL-FOIL~\citep{Fanizzi2018DLFoil} does not support data properties at all.

In this work, we speed up and improve concept learning using a bottom-up approach, dubbed EvoLearner, that is based on random walks and evolutionary algorithms. Instead of refining the top concept $\top$, we start at the instance data by initializing the initial population of the evolutionary algorithm with biased random walks starting from the positive instances in the knowledge graph. We further refine the initial candidates by means of crossover and mutation operations. For data properties, we determine splits to maximize information gain in a way inspired by decision trees~\citep{Quinlan1986Induction}. Our approach significantly outperforms the state-of-the-art approaches CELOE~\cite{Lehmann2009DL-Learner} (top-down), Aleph~\cite{Muggleton1995Inverse} (bottom-up) and SPaCEL~\cite{Tran2012Approach} (hybrid) on 7 of 9 datasets available in the SML-Bench benchmarking framework~\citep{Westphal2019SML-Bench}. Our ablation study shows that both our random-walk initialization and support for data properties significantly contribute to our strong performance. Investigating predictive performance as a function of available runtime shows that EvoLearner outperforms CELOE and SPaCEL at all time points.

In the following, we briefly discuss related work in Section~\ref{sec:related-work} before introducing our EvoLearner in Section~\ref{sec:approach}. Section~\ref{sec:evaluation} evaluates our approach and Section~\ref{sec:discussion} discusses its strength and weaknesses. Section~\ref{sec:conclusion} concludes the paper.
